\documentclass[10pt,letterpaper]{article}

\usepackage[margin=0.5in]{geometry}
\usepackage[T1]{fontenc}
\usepackage{lmodern}
\usepackage{xcolor}
\usepackage{enumitem}
\usepackage{tabularx}
\usepackage[hidelinks]{hyperref}
\usepackage{titlesec}
\usepackage{microtype}

\definecolor{accent}{HTML}{145A78}
\setlength{\parindent}{0pt}
\setlength{\parskip}{0pt}
\setlist[itemize]{leftmargin=1.15em,itemsep=0.8pt,topsep=1pt,parsep=0pt}
\titleformat{\section}{\large\bfseries\color{accent}}{}{0pt}{}[\vspace{-2pt}\titlerule]
\titlespacing*{\section}{0pt}{5pt}{2pt}
\pagestyle{empty}

\newcommand{\role}[4]{%
  \textbf{#1}\hfill #2\\
  \textit{#3}\hfill \textit{#4}\vspace{-3pt}
}

\begin{document}
\fontsize{9}{10.5}\selectfont

\begin{center}
  {\Large\bfseries Piter Garcia}\\[1pt]
  {\large Computer Science Educator and Applied AI Professional}\\[2pt]
  Rochester, NY \enspace|\enspace Available nationwide after December 2026\\
  \href{mailto:pgarcia8@u.rochester.edu}{pgarcia8@u.rochester.edu}
  \enspace|\enspace +1 (646) 288-3300
  \enspace|\enspace \href{https://linkedin.com/in/garciapiterz}{linkedin.com/in/garciapiterz}\\
  \href{https://github.com/pzg8794}{github.com/pzg8794}
  \enspace|\enspace \href{https://garciapiterz.com}{garciapiterz.com}
\end{center}

\section*{Profile}
Computer science educator and applied AI professional preparing to teach,
advise, and mentor students in a K--12 independent-school community. Brings
K--5 student-teaching experience, graduate preparation in inclusive computer
science education, and industry experience building software, machine-learning,
and data systems. Designs technically rigorous learning with multiple ways for
students to participate, communicate, debug, and demonstrate understanding.

\section*{Teaching and Mentoring Experience}
\role{Computer Science Student Teacher}{Spring 2026}
{University of Rochester / Pine Brook Elementary School}{Rochester, NY}
\begin{itemize}
  \item Supported and increasingly led computer-science learning across
    kindergarten through fifth grade using coding, robotics, Sphero, Minecraft
    Education, Tinkercad, Code.org, circuits, digital fluency, and
    student-designed AI projects.
  \item Designed and revised lessons, directions, rubrics, and formative
    assessments using computer-science standards and Universal Design for
    Learning.
  \item Used visual, hands-on, and discussion-based supports to preserve learner
    agency and help students with varied communication, attention, confidence,
    and prior technical experience participate in rigorous work.
\end{itemize}

\role{Tutor and Mentor}{Ongoing}
{Computer Science, Data Science, Statistics, and Mathematics}{Rochester, NY}
\begin{itemize}
  \item Explain technical concepts through examples, diagrams, guided practice,
    and patient debugging, including support for neurodivergent and
    underrepresented learners.
  \item Help learners organize complex assignments, interpret feedback, and
    build confidence while maintaining responsibility for their own work.
\end{itemize}

\role{Graduate Assistant, Quantum and AI Research}{Fall 2025--Present}
{Rochester Institute of Technology, Software Engineering}{Rochester, NY}
\begin{itemize}
  \item Build and document reproducible Python experiments in adaptive
    quantum-network routing, online learning, fairness, and resource allocation.
  \item Translate complex methods and findings into reports, presentations,
    visual explanations, and reusable project materials for varied audiences.
\end{itemize}

\section*{Applied Technical Experience}
\role{Data Solutions Engineer}{Sep 2022--Aug 2024}
{VEDADATA Inc.}{New York, NY}
\begin{itemize}
  \item Built maintainable Python data workflows, validation checks, and
    documentation while collaborating with technical and nontechnical teams.
\end{itemize}

\role{AI Data Solutions Engineer}{Jan 2021--Sep 2022}
{VIOME Inc.}{New York, NY}
\begin{itemize}
  \item Developed and evaluated machine-learning workflows for healthcare data,
    connecting technical decisions with reliability, communication, and user
    impact.
\end{itemize}

\section*{Education}
\textbf{M.S., Teaching Computer Science K--12}, University of Rochester
\hfill Expected December 2026\\
GPA: 3.9+; NSF Robert Noyce Scholar; Advanced Certificate in Leadership in
Disability and Inclusive Practices

\textbf{M.S., Data Science}, Rochester Institute of Technology
\hfill Expected December 2026\\
GPA: 3.9+; applied AI, bioinformatics, neural networks, and reproducible data systems

\textbf{M.S., Computer Science}, Rochester Institute of Technology
\hfill 2015\\
Artificial intelligence, big-data analytics, and computer graphics

\textbf{B.S., Electrical Engineering Technology; B.S., Computer Engineering Technology},
Farmingdale State College \hfill 2012

\section*{Teaching and Technical Strengths}
Computer science and interdisciplinary STEM \enspace|\enspace Inclusive lesson
and assessment design \enspace|\enspace Advising and mentoring \enspace|\enspace
Python, Java, C++, JavaScript, SQL, R, Git/GitHub, Jupyter, and LaTeX

\end{document}
